The settings of the replacement are the row of Section~\ref{sec:settings}, one entry per number of variables of a space, each entry the lengths $d_1$, $d_2$, $d_3$, the pattern set $\Pi$ and, for the slack, the term set $F$ with its offsets $o$; an $\R^q$ space carries one multiplier length. The knobs, as the construction of Sections~\ref{sec:1d}--\ref{sec:nd} defines them:

\begin{enumerate}\setlength\itemsep{0pt}
\item $d_1$, the length of the incomplete moments (block $\bc$, the lower integrals), per direction: the kernel of the certificate has a jump on the diagonal only when $d_1\ge1$.
\item $d_2$, the length of the complete moments (block $\bb$, the full integrals), per direction.
\item $d_3$, the multiplier lengths, per block $\ba,\bb,\bc$ and per direction: the degree in $\theta_k$ of the monomials that the Gram of each block carries; in a direction a space does not own, each basis operator carries its own largest multiplier degree (the rule of the reader, Section~\ref{sec:reader-new}).
\item $\Pi$, the pattern set of a space in two or more variables: \code{uniform} (the pure multiplier and the pure complete moment), \code{separable} ($\{\ba,\bb\}^N$, the mixed operators included) or \code{all} ($\{\ba,\bb,\bc\}^N$).
\item $c$, the per-direction cap: the kernel monomials of every basis operator satisfy (degree in $\theta_k$) $+$ (degree in $s_k$) $\le c_k$; no total-degree cap and no cap table (Section~\ref{sec:cap}).
\item For the slack, $F$ and $o$: the term set and its offsets, and the raises $\mathrm{dD}$, $\mathrm{dw}$ of the reader above the support of the target.
\item For the default stability test, the four caps of the free operator $\mcl Q$ (\code{mult}, \code{int}, \code{out}, \code{in}), which replace the storage row there (Section~\ref{sec:mixed-q}).
\end{enumerate}

\subsection{The reader as implemented}\label{sec:reader-new}

One routine, \code{pf\_row\_storage} in the record (the proposal's reader of Section~\ref{sec:reader} with two measured rules), turns a row into the arguments of \code{poscopvar}: the basis operators are the patterns of $\Pi$ with codes $\ba\to1$, $\bc\to2$ (lower integral) and $\bb\to3$ (upper integral) in a direction where some space carries $\bc$, else $\bb\to4$ (full integral), by Proposition~\ref{prop:span}; each operator gets \code{mult}$_k=\ell(\alpha_k)-1$ and \code{int}$_k=d_{3,\alpha_k,k}-1$ in its owned directions; in a direction its space does not own it carries its own largest \code{int}; a zero length drops the pattern; the per-direction cap goes to \code{poscopvar\_direct} as the field \code{cap}. Validation: the 1-D entry of the light tier reproduces the stock 1-D storage bit for bit on eight library rows and on the ODE--PDE row S1 (same variables, rows and verdicts); the uniform entries reproduce the N-D default of \code{poscopvar} within the dropped total cap (S3: $318\,361$ against $318\,004$ variables, both certified); the separable entry at $d_2=3$ reproduces the S18 certificate.

\subsection{The storage row of each tier}\label{sec:tiers-values}

Table~\ref{tab:tiers} gives the storage row of each tier, one line per entry; every direction of an entry carries the same lengths, and an $\R^q$ space carries the identity basis (\code{len\_R} $=1$) in every tier. The 1-D entries are the present 1-D storages (Table~\ref{tab:1d}), which the row reproduces bit for bit; the 2-D and 3-D entries are fixed by the baselines of Section~\ref{sec:baselines}. Custom carries the storage row of light, and since 11/18/2026 differs from light only in the 1-D degrees of $\mcl Q$ of the direct form (Table~\ref{tab:knobs-q}).

\begin{table}[h]\centering\footnotesize
\setlength{\tabcolsep}{4pt}
\begin{tabular}{llrrrrrllrr}
\toprule
tier & entry & $d_1$ & $d_2$ & $d_{3\ba}$ & $d_{3\bb}$ & $d_{3\bc}$ & $\Pi$ & cap $c$ & cross & Gram rows \\
\midrule
stripped & 1-D & 2 & 2 & 2 & 2 & 2 & \code{all} & $\infty$ & 1 & 10 \\
 & 2-D & 0 & 2 & 1 & 2 & -- & \code{uniform} & $\infty$ & 1 & 17 \\
 & 3-D & 0 & 1 & 1 & 1 & -- & \code{uniform} & $\infty$ & 1 & 2 \\
\midrule
light, custom & 1-D & 2 & 2 & 2 & 2 & 2 & \code{all} & $\infty$ & 1 & 10 \\
 & 2-D & 0 & 2 & 1 & 2 & -- & \code{separable} & $\infty$ & 1 & 25 \\
 & 3-D & 0 & 2 & 1 & 2 & -- & \code{uniform} & $\infty$ & 1 & 65 \\
\midrule
heavy & 1-D & 2 & 2 & 3 & 2 & 2 & \code{all} & $\infty$ & 1 & 11 \\
 & 2-D & 0 & 2 & 2 & 2 & -- & \code{separable} & $\infty$ & 1 & 36 \\
 & 3-D & 0 & 2 & 1 & 2 & -- & \code{uniform} & $\infty$ & 1 & 65 \\
\midrule
veryheavy & 1-D & 4 & 4 & 4 & 4 & 4 & \code{all} & $\infty$ & 1 & 36 \\
 & 2-D & 0 & 3 & 2 & 3 & -- & \code{separable} & $\infty$ & 1 & 121 \\
 & 3-D & 0 & 2 & 1 & 2 & -- & \code{uniform} & $\infty$ & 1 & 65 \\
\midrule
extreme & 1-D & 0 & 0 & 2 & -- & -- & $\ba$ & $\infty$ & 1 & 2 \\
 & 2-D & 0 & 0 & 1 & -- & -- & $(\ba,\ba)$ & $\infty$ & 1 & 1 \\
 & 3-D & 0 & 0 & 1 & -- & -- & $(\ba,\ba,\ba)$ & $\infty$ & 1 & 1 \\
\bottomrule
\end{tabular}
\caption{The storage rows of the proposed tiers. Lengths per direction ($0$: the block is absent); $\Pi$ the pattern set; $c$ the per-direction cap of Section~\ref{sec:cap} ($\infty$: none); cross the cross length \code{len\_cross}, the length in a registry direction the space does not own, read on a mixed registry only ($1$: none); the last column the rows of the plain Gram per state component, $\sum_{\alpha\in\Pi}\prod_k\ell_k(\alpha_k)\,d_{3,\alpha_k,k}$ under the cap. The storage carries no term in any tier; the terms of the slack, the same set in every tier, are listed in Section~\ref{sec:tier-terms}.}\label{tab:tiers}
\end{table}

The lengths of every entry are those of the present tier (Tables~\ref{tab:1d}--\ref{tab:nd}); the row changes three things. The pattern set of a 2-D entry is \code{separable} in light, heavy and veryheavy, where the present stock 2-D storage is \code{uniform}: at the light lengths the separable row costs the stock size on the 2-D library and the 2-D ladder (Table~\ref{tab:bl-lib}; within $0.6\%$ in variables) and certifies the mixed suite where the stock 2-D storage fails S3 and S8 (Table~\ref{tab:bl-mixed}), and at the veryheavy lengths it is the one storage that certifies S18, at both margins (Section~\ref{sec:mixed-patterns}). The total-degree caps are gone: the uncapped row at the light lengths reproduces the stock size in two variables, because the storage Gram is a small part of the SDP once the slack is sized from it, and in three variables adds $0.4\%$ to the variables of the 3-D heat benchmark (Table~\ref{tab:bl-fin2}); the per-direction cap is kept as the knob for the entries that need it, with the measured trade at the light lengths that $c=1$ returns an \code{all} entry with $d_1=1$ to the stock size at $6\%$ of the fin2 gap (Table~\ref{tab:bl-cap}). The cross length is $1$ in every entry, where the present construction carries the rule of the reader ($2$ at the light lengths, $3$ at the veryheavy ones), and on the mixed suite it certifies every row the rule certifies, at $0.2$ to $1\%$ fewer variables (Section~\ref{sec:bl-cross}).

Two rows measured in this report are not tiers but options a caller sets on a row: \code{'all'} with $d_1\ge1$ in two variables (Rlight: the storage that contains the PDE norm, Section~\ref{sec:mixed-d1}; $2\times$ the variables on the mixed gains and $1.8$--$5\times$ on the 2-D library, with the mixed gains tighter by $1.6$--$3\times$, and at $c=1$ the stock size), and the uniform row (Rdef, the present \code{poscopvar} default, within $0.4\%$ of the stock size everywhere and $12\%$ looser on PDE~48). On PDEs with polynomial coefficients the separable row again costs the stock size and changes nothing on single domains, and on the mixed cascade with variable coefficients it reaches $0.977$ of the threshold where the stock light storage reaches $0.736$ (Section~\ref{sec:bl-inhom}). The light 2-D entry was chosen among three stock-sized rows that replicate the library and the mixed suite alike (Table~\ref{tab:bl-light}): the separable row is the one contained in the heavy and veryheavy entries, so the tiers nest; the capped all-pattern row is the option for a caller who needs the incomplete moments in two variables at the light cost.

Three measured limits of the proposal. The heavy 2-D entry reproduces the size of the stock heavy 2-D storage (fin2: $760\,485$ against $759\,358$ variables) and not its fin2 gap ($19$ against $3.7$ in units of $10^{-5}$; the uniform row at the same lengths $15$): the stock cap tables admit a storage the row at these lengths does not express (Section~\ref{sec:cannot}), and no row of this report reaches that gap below the veryheavy lengths ($2.2$ at $12\times$ the size). Gaps on fin2 within a factor $1.3$ are inside the solver's variation: two nested pairs came out inverted (the uniform row $48$ against stock $47$ at light; the separable row $19$ against the uniform $15$ at heavy), so a gap ranks rows only beyond that factor. In three variables the stability executive at light is infeasible on the 3-D heat at $0.9\kappa^*$ under both the stock storage and the row, at $\mathrm{dw}=0$ with no raise allowed and the slack bases bounded by the degree of $\mcl T$ (\code{cap\_delta} $0$): the 3-D entry of the row is fixed by the N-D default it reproduces and by the heat benchmark record (Section~\ref{sec:open}), not by an executive baseline. Without that bound the light slack in three variables has $20.9$M decision variables (seven PSD blocks of $1\,728$ rows, $m=126\,681$) against $908$k with it, so the bound is what makes the 3-D light tier affordable, and the raise it forbids is the open item, not the storage.


\subsection{The slack settings of each tier}\label{sec:knobs-slack}

The slack of every LPI is sized by the reader from its target: lift length $D_{\min}+\mathrm{dD}$, weight length $w+\mathrm{dw}$ with the floor \code{wmin} when the target is the storage itself, and the term set $F$ with the offsets $o$ of the option \code{psatz} (Section~\ref{sec:tier-terms}); the degree loop of the executives raises the weight after UNKNOWN or INFEASIBLE, from the tier's \code{dw\_feas} (stability, well-posedness), \code{dw\_P} (coercive gains, norms, synthesis) or \code{dw\_Q} (the non-coercive forms), at most \code{max\_raise\_feas} or \code{max\_raise} times, the slack bases bounded per direction by the largest degree of $\mcl T$ plus \code{cap\_delta}. Table~\ref{tab:knobs-slack} gives \code{st.lpi.ineq} and the free-operator degrees \code{lpi.var.deg} of each tier, Table~\ref{tab:knobs-raise} the raise knobs of \code{st.sop}. Every value is that of the present tier files, which the measurements of Sections~\ref{sec:fix} and \ref{sec:program} leave in place, with three changes: the 2-D entry of \code{lpi.var.deg} is raised to $2$ in heavy and veryheavy for the direct default test (Section~\ref{sec:mixed-q}), where degree $2$ certified S17 at $0.5r_2^*$ with $\mathrm{dw}=1$ and degree $1$ did not (Table~\ref{tab:qdirect}); \code{wcross} is $0$; and the 2-D term set is \code{'faces1'} (Section~\ref{sec:slack-settings}; \gapsref).

\begin{table}[h]\centering\scriptsize
\setlength{\tabcolsep}{4pt}
\begin{tabular}{lcccclll}
\toprule
tier & \code{dD} & \code{dw} & \code{wmin} & \code{wcross} & \code{psatz} ($F$) & \code{psatz\_offset} ($o$) & \code{lpi.var.deg} ($\mcl Q$) \\
\midrule
stripped & 1 & 1 & 0 & 0 & \code{\{'auto','faces1','auto'\}} & \code{[]} (the rule) & $[2\,2\,2]$, 1, 1 \\
light, custom & 1 & 1 & $[2\,2\,0]$ & 0 & \code{\{'auto','faces1','auto'\}} & \code{[]} & $[2\,2\,2]$, 1, 1 \\
heavy & 1 & 1 & $[2\,2\,0]$ & 0 & \code{\{'auto','faces1','auto'\}} & \code{[]} & $[4\,2\,2]$, 2, 1 \\
veryheavy & 1 & 1 & $[2\,2\,0]$ & 0 & \code{\{'auto','faces1','auto'\}} & \code{[]} & $[6\,6\,6]$, 2, 1 \\
extreme & 1 & 1 & 0 & 0 & \code{\{'auto','faces1','auto'\}} & \code{[]} & $[2\,2\,0]$, 1, 1 \\
\bottomrule
\end{tabular}
\caption{The slack settings \code{st.lpi.ineq} of each tier and the free-operator degrees \code{lpi.var.deg} (entries for one, two and three or more variables). \code{psatz} is a cell per number of registry variables: \code{'auto'} in one variable, the box quadratic; \code{'faces1'} in two, the four faces at offset $1$ whatever the weight degree; \code{'auto'} in three or more, the $2N$ faces at the rule of Section~\ref{sec:tier-terms}, offset $1$ when every weight degree is at least $2$, else $0$. \code{wcross} $0$: no weight in a registry direction a space does not own. $\varepsilon_{\mathrm{pos}}=10^{-4}$, $\varepsilon_{\mathrm{neg}}=0$ in every tier.}\label{tab:knobs-slack}
\end{table}

\begin{table}[h]\centering\scriptsize
\setlength{\tabcolsep}{4pt}
\begin{tabular}{lcccccc}
\toprule
tier & \code{dw\_feas} & \code{dw\_P} & \code{dw\_Q} & \code{max\_raise\_feas} & \code{max\_raise} & \code{cap\_delta} \\
\midrule
stripped & $(0,1,1)$ & 1 & 1 & $(1,1,0)$ & 0 & $(\cdot,\cdot,0)$ \\
light, custom & $(0,1,1)$ & $(2,2,1)$ & 1 & $(2,1,0)$ & $(1,1,0)$ & $(\cdot,\cdot,0)$ \\
heavy & $(1,2,2)$ & 2 & 1 & $(2,2,0)$ & $(1,1,0)$ & $(\cdot,\cdot,1)$ \\
veryheavy & $(1,2,2)$ & 2 & 1 & $(3,3,0)$ & $(2,2,0)$ & $(\cdot,\cdot,2)$ \\
extreme & $(0,1,1)$ & 1 & 1 & $(1,1,0)$ & 0 & $(\cdot,\cdot,0)$ \\
\bottomrule
\end{tabular}
\caption{The raise knobs \code{st.sop} of each tier: the weight raise by LPI, the number of raises the degree loop allows, and the bound on the slack bases ($\cdot$: none); triples are the entries for one, two and three or more variables.}\label{tab:knobs-raise}
\end{table}

\paragraph{Implemented 11/17/2026.} The six tier files carry the values of Tables~\ref{tab:tiers}, \ref{tab:knobs-slack} and \ref{tab:knobs-raise} in the layout of Section~\ref{sec:tier-layout} (work package 5, \file{scratchpad/wp5/WP5_RESULTS.md}; \file{DECISIONS.md}, item D73), \code{lpi.var.deg} as above: \code{\{[2 2 2],1,1\}} in light, custom and stripped, \code{\{[4 2 2],2,1\}} in heavy, \code{\{[6 6 6],2,1\}} in veryheavy, \code{\{[2 2 0],1,1\}} in extreme. Verification through the executives at light, MOSEK. The 1-D rows (PDE~1, 2, 3, 5, 13, 17, 21 and S1) are identical in variables, rows and verdict to the repository before the change in \code{'P'} and to work package 4b in \code{'Q'}. The 2-D rows, \code{'P'} / \code{'Q'}, all certified: PDE~33 $26\,545$ / $52\,660$; S3 $84\,205$ / $86\,078$; S8 $244\,521$ / $101\,424$; S26 $246\,681$ / $69\,915$. The stock sizes of Section~\ref{sec:baselines} are PDE~33 $103\,744$, S3 $508\,228$ and S8 $1\,336\,720$ (both UNKNOWN) and S26 $336\,926$ (Tables~\ref{tab:bl-lib} and \ref{tab:bl-mixed}). S17 at $0.5r_2^*$: \code{'Q'} certified at $98\,410$, where \code{'auto'} with \code{wcross} $0$ was UNKNOWN after one raise at $676\,410$. S18 at $0.5\tau^*$: \code{'P'} infeasible after one raise at $428\,185$ (stock $508\,228$, infeasible), \code{'Q'} certified after one raise at $374\,827$, where \code{'auto'} certified with no raise at $201\,050$. The 3-D heat at $0.9\kappa^*$: \code{'P'} infeasible at $911\,425$ variables ($907\,729$ recorded; $+0.4\%$). The $H_\infty$ gain of the LRM case of the literature set: $0.0459441$ at $203\,522$ variables, between the stock and row records ($0.0459440$ and $0.0459459$, $210\,006$ variables).


\subsection{The terms of each tier}\label{sec:tier-terms}

Every tier carries the same term set, and carries it in the slack only; the tiers differ in the degrees of the terms, not in the set. The storage of every tier carries no term ($F=$ \code{none} in every entry of Table~\ref{tab:tiers}), as the present tiers (\code{LF\_use\_psatz} $0$). The slack of every inequality, the two inequalities of the direct default test included, is certified by the reader with the plain term and the weights of the rule \code{'auto'} of \code{get\_lift\_degs}, Table~\ref{tab:tier-terms}; since 11/17/2026 the 2-D entry is \code{'faces1'}, the same faces at offset $1$ whatever the weight degree (Section~\ref{sec:slack-settings}).

\begin{table}[h]\centering\scriptsize
\setlength{\tabcolsep}{4pt}
\begin{tabular}{llll}
\toprule
variables & weights $g$ of the slack & offset $o_g$ & multiplier lengths \\
\midrule
1 & plain; the box product $(s-a)(b-s)$ & $0$; $1$ & $w+1$; $w$ \\
2 & plain; the four faces $s_k-a_k$, $b_k-s_k$ & $0$; $1$ (\code{'faces1'}) & $w_k+1$; $w_k$ \\
3 & plain; the six faces & $0$; $1$ if $\min_k w_k\ge2$, else $0$ & $w_k+1$; $w_k+1-o$ \\
\bottomrule
\end{tabular}
\caption{The terms of the slack in every tier (\code{'auto'}; in two variables \code{'faces1'}): the weights, their offsets (the number of degrees by which the multiplier vector of a weighted term is shorter than that of the unweighted term) and the length of the multiplier monomial vector of each term in direction $k$; the moment lengths of every term are $D_k+1$. In one variable the plain term and the box product at offset $1$ are the Markov--Luk\'acs pair.}\label{tab:tier-terms}
\end{table}

The lengths come from the target. With $D_{\max,k}$, $M_k$ and $J_k$ the degrees of the support of the target in direction $k$ (its lower kernels, its multiplier, its couplings with $\R^q$) and $w_R$ the weight degree of an $\R^q$ space,
\[
D_k=D_{\min,k}+\mathrm{dD},\qquad w_k=\max\Bigl\{\Bigl\lceil\tfrac{D_{\max,k}-D_k-1}{2}\Bigr\rceil,\ \Bigl\lceil\tfrac{M_k}{2}\Bigr\rceil,\ J_k-D_k-w_R-1,\ 0\Bigr\}+\mathrm{dw},
\]
raised to \code{wmin} when the target is the storage itself. At $w_k=1$ the faces stand at the degree of the plain term and reach total degree $2w+1$, the half step that certifies S18 at half its margin (Section~\ref{sec:t2}); at $w_k\ge2$ they stand one degree below it. The tier enters through $\mathrm{dw}$ (\code{dw\_feas} for stability and well-posedness, \code{dw\_P} for the coercive gains, norms and synthesis, \code{dw\_Q} for the non-coercive forms), through \code{wmin}, and through the raises of Table~\ref{tab:knobs-slack}: in two variables light sizes the stability slack at $w=$ reach $+1$ and the gains at reach $+2$, heavy both at reach $+2$ with one more raise allowed, veryheavy with two more. The box set and the corners enter no tier: on the mixed suite neither certifies a case the faces do not (Table~\ref{tab:t2}), at $1.2\times$ and $1.8\times$ the variables of the scalar-offset faces on S3. The standing of this choice in the theory (Section~\ref{sec:terms}): in one variable the tier carries the product, the degree-exact form of Lemma~\ref{lem:ml}; in $N\ge2$ the faces are complete for scalar polynomials by Putinar and the products of faces by Schm\"udgen, no set is proven complete for the kernels of a PI operator, and the faces are the measured choice. The product sets the theory names, the preordering of the box quadratics, the faces with the face-by-box products across directions, and the full preordering of the faces, are measured on the mixed suite in Section~\ref{sec:t3}, and the alternative to the odd-degree faces, one more weight degree with even-degree weights, in Section~\ref{sec:t4}: on S18 it certifies both margins at $1.9\times$ the variables of the full-length faces at the lower degree, on S17 it certifies less, and in three variables it certifies the heat benchmark at $0.985\kappa^*$ at $5.1\times$ the variables and, on $64$ threads, $24\times$ the time of the faces (Section~\ref{sec:t4}).
